\documentclass[DIN, pagenumber=false, fontsize=11pt, parskip=half]{scrartcl}
\usepackage[UTF8]{ctex}  % 中文支持
% \usepackage{ngerman}
\usepackage[utf8]{inputenc}
\usepackage[T1]{fontenc}
\usepackage{textcomp}
\usepackage{xcolor} % Required for custom red

\usepackage{geometry} % Custom margins

\usepackage{tocloft} % Modifying the table of contents

% for matlab code
% bw = blackwhite - optimized for print, otherwise source is colored
\usepackage[framed,numbered,bw]{mcode}
\usepackage[most]{tcolorbox}
% for other code
%\usepackage{listings}

\setlength{\parindent}{0em}

% set section in CM
\setkomafont{section}{\normalfont\bfseries\Large}

\newcommand{\mytitle}[1]{{\noindent\Large\textbf{#1}}}
\newcommand{\mysection}[1]{\textbf{\section*{#1}}}
\newcommand{\mysubsection}[2]{\romannumeral #1) #2}

\newtcolorbox{parchmentquote}{
    enhanced,
    colback=orange!5!white,   % 浅米黄色背景
    colframe=orange!50!black,
    boxrule=0.5pt,
    arc=2pt,                  % 细微圆角
    boxsep=5pt,
    before skip=10pt,
    after skip=10pt,
    fontupper=\itshape,
    borderline west={2pt}{0pt}{orange!80!black} % 加厚的左装饰线
}
%===================================
\begin{document}

\noindent\textbf{课下练习} \hfill \textbf{online course}\\
primary school \hfill Liu\\

\mytitle{分数复习 \hfill \today}
\newline
\begin{parchmentquote}
    对分数的概念、基本性质与计算进行复习。\newline
    最大公约数与最小公倍数短除法要掌握，在分数运算中约分、通分有广泛应用。
\end{parchmentquote}


%===================================
\section{分数概念}

\subsection{分数读写}
\subsubsection{读出下面分数}

\begin{enumerate}
    \item $\frac{1}{2}$
    \item $\frac{3}{2}$
    \item $5\frac{6}{11}$
\end{enumerate}
\subsubsection{写出下面分数}

\begin{enumerate}
    \item 九十五分之七
    \item 七分之七
    \item 三又九分之五
\end{enumerate}
\subsection{分数概念的理解}
\subsubsection{假分数转成带分数}
\begin{enumerate}
    \item $\frac{5}{3}$
    \item $\frac{136}{11}$
\end{enumerate}
\subsubsection{带分数转成假分数}
\begin{enumerate}
    \item $4\frac{2}{3}$
    \item $12\frac{7}{18}$
\end{enumerate}
\subsubsection{分数单位}
\begin{enumerate}
    \item $\frac{2}{3}$的分数单位
    \item 以$\frac{1}{18}$为分数单位的真分数有哪些
\end{enumerate}
\section{分数运算}
\begin{enumerate}
    \item $\frac{2}{3} + \frac{5}{6} = $
    \item $\frac{7}{14} + \frac{5}{21} = $
    \item $1.5 \times \frac{2}{5}  = $
    \item $\frac{9}{25} \times \frac{20}{21}  = $
    \item $\frac{12}{49} \div \frac{21}{56}  = $
\end{enumerate}

\end{document}